\documentclass[11pt]{letter}

\usepackage[margin=0.5in]{geometry}
\usepackage{hyperref}

\begin{document}

\begin{letter}{}

\opening{Dear Hiring Team,}

I am writing to apply for the Machine Learning Engineer, New Grad position at Quora. With a 2026 Master of Science in Computer Science from the University of Calgary and hands-on experience building Python/C++ machine learning systems, LLM-enabled automation workflows, NLP-related backend services, and model evaluation pipelines, I am excited by the opportunity to contribute to consumer-facing AI systems at Quora.

What stands out to me is the role's focus on building AI-powered product experiences, not only training models in isolation. The combination of LLM and generative AI features, prompt and workflow iteration, offline evaluation, user engagement and trust metrics, and close product collaboration fits the kind of work I want to grow into. I am especially interested in Quora because its products connect AI systems with real information needs at very large scale.

My graduate research strengthened the core habits I would bring to this role: careful experimentation, robust data processing, and honest evaluation. In my DEM super-resolution project, I trained and evaluated PyTorch models, built data pipelines with Rasterio and the Google Earth Engine API, and used quantitative metrics such as RMSE, slope, terrain ruggedness, topographic position, and wavelet-based errors to evaluate model behavior. I also used geographically separated validation to reduce data leakage and get a more realistic view of model generalization.

I also bring practical software and NLP-related experience. As a back-end developer at Asr Gooyesh Pardaz, I built Django REST APIs for NLP, text-to-speech, and speech-to-text services backed by PostgreSQL, implemented authentication flows, and contributed to a chatbot system using GPT-style responses and FAQ data. More recently, I have been building an agentic job-search automation platform using LLM workflows, n8n, OpenClaw, Docker, Google Sheets, Google Drive, SearXNG, and Telegram API integrations. That project involves prompt-driven analysis, structured outputs, workflow automation, and practical decision support.

My research software work with BigGeo, Mitacs, and the University of Calgary also gave me experience taking open-ended technical problems through design, implementation, debugging, performance measurement, and documentation. I improved one C++ processing workflow from 233 seconds to 26 seconds for a 10-tensor workload and published my M.Sc. research as first author. I have not yet operated production-scale ML serving systems at Quora's level, but I bring strong Python/C++ fundamentals, ML experimentation experience, LLM workflow practice, persistence, and a genuine interest in building AI products that are useful, reliable, and carefully evaluated.

Thank you for considering my application. I would welcome the opportunity to discuss how my machine learning, LLM workflow, data systems, NLP/backend, and research software experience can contribute to Quora's AI product engineering work.

\closing{Sincerely,

Amir Mirzai Golpayegani

\href{mailto:amirgolpaa24@gmail.com}{amirgolpaa24@gmail.com}}

\end{letter}

\end{document}
